% GENERATED FILE. Edit canonical lessons, then run npm run generate:books.
% canonical-source-hash: fnv1a64:b9d36167d3b70c7e
% canonical-lessons: MW-C234-evaliyo, MW-C234-lathi, MW-C234-charno, MW-C234-hakno, MW-C234-sambhalno, MW-R234-review-229-234

\chapter{The Herd}
\label{ch:mw-the-herd}
\hlchaptermodality{234}

\begin{chapteropening}
\textbf{By the end of this chapter:} \emph{I can name the shepherd and say how the animals are driven and grazed.}

Complete the last lesson of chapter 234: I can name the shepherd and say how the animals are driven and grazed.
\end{chapteropening}

\section[\texorpdfstring{\mw{एवाळियो} (evāḷiyo)}{एवाळियो (evāḷiyo)}]{\mw{एवाळियो} (evāḷiyo) — a shepherd}

\label{lesson:MW-C234-evaliyo}



\begin{quote}
Before the new one: say the Marwadi for to tighten, then the Marwadi for a small bell.
\end{quote}

\subsection*{You'll want to know: \mw{एवाळियो}}
\textbf{\mw{एवाळियो}} — \emph{evāḷiyo} — \textquotedblleft{}a shepherd\textquotedblright{}.

\begin{quote}
\mw{एवाळियो} \mw{जंगळ} \mw{मांय} \mw{जावै} \mw{है।}
\end{quote}

\begin{quote}
\emph{evāḷiyo jaṅgaḷ mā̃y jāvai hai.}
\end{quote}

\begin{quote}
The shepherd goes out into the scrubland.
\end{quote}

A masculine noun.

\begin{cousinweb}
\textbf{\mw{एवाळियो}} (\emph{evāḷiyo}) is inherited from Sanskrit \emph{avipāla}, \textquotedblleft{}a shepherd\textquotedblright{}, built from \emph{avi}, \textquotedblleft{}sheep\textquotedblright{}, and \emph{pāla}, \textquotedblleft{}a keeper\textquotedblright{}; CDIAL lists Marwari \emph{evāḷiyo} (CDIAL 893).
\end{cousinweb}

\begin{grammarlens}[title={What you've built}]
One more word for this chapter.
\end{grammarlens}

\subsection*{Your turn}
\begin{itemize}
\raggedright
  \item \emph{Say it:} \emph{evāḷiyo}
  \item \emph{Say it:} \emph{evāḷiyo}, once more
  \item \emph{Say it:} \emph{evāḷiyo jaṅgaḷ mā̃y jāvai hai.} — The shepherd goes out into the scrubland.
  \item \emph{Recall:} say the Marwadi for to tighten, then the Marwadi for a small bell, then say \emph{evāḷiyo} again
\end{itemize}

\begin{tcolorbox}[breakable,colback=teal!4,colframe=teal!35!black,title={Before you move on}]
What does \mw{एवाळियो} mean? (\textquotedblleft{}A shepherd\textquotedblright{}.) Say it once more.
\end{tcolorbox}
\section[\texorpdfstring{\mw{लाठी} (lāṭhī)}{लाठी (lāṭhī)}]{\mw{लाठी} (lāṭhī) — a long stick}

\label{lesson:MW-C234-lathi}



\begin{quote}
Before the new one: say the Marwadi for a small bell, then the Marwadi for a shepherd.
\end{quote}

\subsection*{You'll want to know: \mw{लाठी}}
\textbf{\mw{लाठी}} — \emph{lāṭhī} — \textquotedblleft{}a long stick\textquotedblright{}.

\begin{quote}
\mw{एवाळियो} \mw{लाठी} \mw{राखै} \mw{है।}
\end{quote}

\begin{quote}
\emph{evāḷiyo lāṭhī rākhai hai.}
\end{quote}

\begin{quote}
The shepherd carries a long stick.
\end{quote}

A feminine noun.

\begin{cousinweb}
\textbf{\mw{लाठी}} (\emph{lāṭhī}) is inherited from a reconstructed \emph{*laṣṭi}, \textquotedblleft{}a stick\textquotedblright{}, beside Sanskrit \emph{yaṣṭi}; CDIAL lists Marwari \emph{lāṭhī} (CDIAL 10991).
\end{cousinweb}

\begin{grammarlens}[title={What you've built}]
One more word for this chapter.
\end{grammarlens}

\subsection*{Your turn}
\begin{itemize}
\raggedright
  \item \emph{Say it:} \emph{lāṭhī}
  \item \emph{Say it:} \emph{lāṭhī}, once more
  \item \emph{Say it:} \emph{evāḷiyo lāṭhī rākhai hai.} — The shepherd carries a long stick.
  \item \emph{Recall:} say the Marwadi for a small bell, then the Marwadi for a shepherd, then say \emph{lāṭhī} again
\end{itemize}

\begin{tcolorbox}[breakable,colback=teal!4,colframe=teal!35!black,title={Before you move on}]
What does \mw{लाठी} mean? (\textquotedblleft{}A long stick\textquotedblright{}.) Say it once more.
\end{tcolorbox}
\section[\texorpdfstring{\mw{चरणो} (charṇo)}{चरणो (charṇo)}]{\mw{चरणो} (charṇo) — to graze}

\label{lesson:MW-C234-charno}



\begin{quote}
Before the new one: say the Marwadi for a shepherd, then the Marwadi for a long stick.
\end{quote}

\subsection*{You'll want to know: \mw{चरणो}}
\textbf{\mw{चरणो}} — \emph{charṇo} — \textquotedblleft{}to graze\textquotedblright{}.

\begin{quote}
\mw{भेड़} \mw{खेत} \mw{मांय} \mw{चरै} \mw{है।}
\end{quote}

\begin{quote}
\emph{bheṛ khet mā̃y charai hai.}
\end{quote}

\begin{quote}
The sheep graze in the field.
\end{quote}

Past (it …): \mw{वो} \mw{ऊंट} \mw{चर्यो}, \emph{vo ū̃ṭ charyo}.

\begin{cousinweb}
\textbf{\mw{चरणो}} (\emph{charṇo}) is inherited from Sanskrit \emph{carati}, \textquotedblleft{}moves, grazes\textquotedblright{}; CDIAL lists the Old Marwari verb \emph{caraï}, \textquotedblleft{}grazes\textquotedblright{} (CDIAL 4686).
\end{cousinweb}

\begin{grammarlens}[title={What you've built}]
One more word for this chapter.
\end{grammarlens}

\subsection*{Your turn}
\begin{itemize}
\raggedright
  \item \emph{Say it:} \emph{charṇo}
  \item \emph{Say it:} \emph{charṇo}, once more
  \item \emph{Say it:} \emph{bheṛ khet mā̃y charai hai.} — The sheep graze in the field.
  \item \emph{Recall:} say the Marwadi for a shepherd, then the Marwadi for a long stick, then say \emph{charṇo} again
\end{itemize}

\begin{tcolorbox}[breakable,colback=teal!4,colframe=teal!35!black,title={Before you move on}]
What does \mw{चरणो} mean? (\textquotedblleft{}To graze\textquotedblright{}.) Say it once more.
\end{tcolorbox}
\section[\texorpdfstring{\mw{हांकणो} (hā̃kṇo)}{हांकणो (hā̃kṇo)}]{\mw{हांकणो} (hā̃kṇo) — to drive (animals or a cart)}

\label{lesson:MW-C234-hakno}



\begin{quote}
Before the new one: say the Marwadi for a long stick, then the Marwadi for to graze.
\end{quote}

\subsection*{You'll want to know: \mw{हांकणो}}
\textbf{\mw{हांकणो}} — \emph{hā̃kṇo} — \textquotedblleft{}to drive (animals or a cart)\textquotedblright{}.

\begin{quote}
\mw{हूं} \mw{काल} \mw{ऊंट} \mw{हांकसूं।}
\end{quote}

\begin{quote}
\emph{hū̃ kāl ū̃ṭ hā̃ksū̃.}
\end{quote}

\begin{quote}
Tomorrow I'll drive the camel.
\end{quote}

Past (I …): \mw{म्हैं} \mw{ऊंट} \mw{हांक्यो}, \emph{mhaĩ ū̃ṭ hā̃kyo}.

\begin{cousinweb}
\textbf{\mw{हांकणो}} (\emph{hā̃kṇo}) is inherited from Sanskrit \emph{hakkayati}, \textquotedblleft{}calls to\textquotedblright{}; CDIAL lists the Old Marwari verb, \textquotedblleft{}drives\textquotedblright{}, under it (CDIAL 13939): a herder drives the animals by calling to them.
\end{cousinweb}

\begin{grammarlens}[title={What you've built}]
One more word for this chapter.
\end{grammarlens}

\subsection*{Your turn}
\begin{itemize}
\raggedright
  \item \emph{Say it:} \emph{hā̃kṇo}
  \item \emph{Say it:} \emph{hā̃kṇo}, once more
  \item \emph{Say it:} \emph{hū̃ kāl ū̃ṭ hā̃ksū̃.} — Tomorrow I'll drive the camel.
  \item \emph{Recall:} say the Marwadi for a long stick, then the Marwadi for to graze, then say \emph{hā̃kṇo} again
\end{itemize}

\begin{tcolorbox}[breakable,colback=teal!4,colframe=teal!35!black,title={Before you move on}]
What does \mw{हांकणो} mean? (\textquotedblleft{}To drive (animals or a cart)\textquotedblright{}.) Say it once more.
\end{tcolorbox}
\section[\texorpdfstring{\mw{संभाळणो} (sambhāḷṇo)}{संभाळणो (sambhāḷṇo)}]{\mw{संभाळणो} (sambhāḷṇo) — to look after; to take care of}

\label{lesson:MW-C234-sambhalno}



\begin{quote}
Before the new one: say the Marwadi for to graze, then the Marwadi for to drive (animals or a cart).
\end{quote}

\subsection*{You'll want to know: \mw{संभाळणो}}
\textbf{\mw{संभाळणो}} — \emph{sambhāḷṇo} — \textquotedblleft{}to look after; to take care of\textquotedblright{}.

\begin{quote}
\mw{हूं} \mw{गाय} \mw{संभाळूं} \mw{हूं।}
\end{quote}

\begin{quote}
\emph{hū̃ gāy sambhāḷū̃ hū̃.}
\end{quote}

\begin{quote}
I look after the cow.
\end{quote}

Past (I …): \mw{म्हैं} \mw{गाय} \mw{संभाळी}, \emph{mhaĩ gāy sambhāḷī}.

\begin{cousinweb}
\textbf{\mw{संभाळणो}} (\emph{sambhāḷṇo}) is inherited from Sanskrit \emph{sambharati}, \textquotedblleft{}brings together\textquotedblright{}; CDIAL lists the Old Marwari verb \emph{samālaï}, \textquotedblleft{}takes care of\textquotedblright{} (CDIAL 12961).
\end{cousinweb}

\begin{grammarlens}[title={What you've built}]
One more word for this chapter.
\end{grammarlens}

\subsection*{Your turn}
\begin{itemize}
\raggedright
  \item \emph{Say it:} \emph{sambhāḷṇo}
  \item \emph{Say it:} \emph{sambhāḷṇo}, once more
  \item \emph{Say it:} \emph{hū̃ gāy sambhāḷū̃ hū̃.} — I look after the cow.
  \item \emph{Recall:} say the Marwadi for to graze, then the Marwadi for to drive (animals or a cart), then say \emph{sambhāḷṇo} again
\end{itemize}

\begin{tcolorbox}[breakable,colback=teal!4,colframe=teal!35!black,title={Before you move on}]
What does \mw{संभाळणो} mean? (\textquotedblleft{}To look after; to take care of\textquotedblright{}.) Say it once more.
\end{tcolorbox}
\section[(dialogue)]{Looking back: from Light and White to The Herd}

\label{lesson:MW-R234-review-229-234}



\begin{quote}
Say the Marwadi for daylight and bright.
\end{quote}

\subsection*{Your turn}
Say these aloud:

\begin{itemize}
\raggedright
  \item daylight, bright, fine, white, to shine
  \item to be scorching, to go down (of the sun), the earth, swallow-wort, a thorn
  \item crooked, straight, standing, round, loose
  \item elderly, soil, smooth, a pebble, a jujube bush
  \item a camel saddle, a camel's nose-peg, trappings (of a camel or a horse), to tighten, a small bell
  \item a shepherd, a long stick, to graze, to drive (animals or a cart), to look after
\end{itemize}

\begin{tcolorbox}[breakable,colback=teal!4,colframe=teal!35!black,title={Before you move on}]
Daylight? (\textbf{\mw{उजास}}, \emph{ujās}.) Bright? (\textbf{\mw{ऊजळो}}, \emph{ūjḷo}.) Fine? (\textbf{\mw{झीणो}}, \emph{jhīṇo}.) White? (\textbf{\mw{धोळो}}, \emph{dhoḷo}.) To shine? (\textbf{\mw{चमकणो}}, \emph{chamakṇo}.) To be scorching? (\textbf{\mw{बळणो}}, \emph{baḷṇo}.) To go down (of the sun)? (\textbf{\mw{ढळणो}}, \emph{ḍhaḷṇo}.) The earth? (\textbf{\mw{धरती}}, \emph{dhartī}.) Swallow-wort? (\textbf{\mw{आक}}, \emph{āk}.) A thorn? (\textbf{\mw{कांटो}}, \emph{kā̃ṭo}.) Crooked? (\textbf{\mw{बांको}}, \emph{bā̃ko}.) Straight? (\textbf{\mw{सीधो}}, \emph{sīdho}.) Standing? (\textbf{\mw{ऊभो}}, \emph{ūbho}.) Round? (\textbf{\mw{गोळ}}, \emph{goḷ}.) Loose? (\textbf{\mw{ढीलो}}, \emph{ḍhīlo}.) Elderly? (\textbf{\mw{बूढो}}, \emph{būḍho}.) Soil? (\textbf{\mw{माटी}}, \emph{māṭī}.) Smooth? (\textbf{\mw{चिकणो}}, \emph{chikṇo}.) A pebble? (\textbf{\mw{कांकरो}}, \emph{kā̃kro}.) A jujube bush? (\textbf{\mw{बोरड़ी}}, \emph{borṛī}.) A camel saddle? (\textbf{\mw{पलाण}}, \emph{palāṇ}.) A camel's nose-peg? (\textbf{\mw{नकेल}}, \emph{nakel}.) Trappings (of a camel or a horse)? (\textbf{\mw{साज}}, \emph{sāj}.) To tighten? (\textbf{\mw{कसणो}}, \emph{kasṇo}.) A small bell? (\textbf{\mw{घूघरो}}, \emph{ghūghro}.) A shepherd? (\textbf{\mw{एवाळियो}}, \emph{evāḷiyo}.) A long stick? (\textbf{\mw{लाठी}}, \emph{lāṭhī}.) To graze? (\textbf{\mw{चरणो}}, \emph{charṇo}.) To drive (animals or a cart)? (\textbf{\mw{हांकणो}}, \emph{hā̃kṇo}.) To look after? (\textbf{\mw{संभाळणो}}, \emph{sambhāḷṇo}.)
\end{tcolorbox}
